\chapter{矩阵函数的微分}
\label{sec:multivariable-matrix-perturbation}
\label{sec:multivariable-matrix-differentiation}
\label{chap:matrix-differentiation}

机器学习中的参数常排成矩阵，例如线性层的权重或协方差矩阵。损失是这些矩阵的标量函数，求逆、相乘等运算又产生矩阵输出。本章要回答：以矩阵为自变量时，微分是什么形状的对象；标量输出的梯度怎样读出；乘法、求逆、行列式与对数行列式这些常用运算的微分公式是什么，在哪里成立。

原则上可以把矩阵拉直成长向量，再用第~\ref{chap:multivariate-differentiation}章的结论。但保留矩阵形状，乘法次序与迹的结构更清楚，结果也能直接写成矩阵表达式。矩阵空间是有限维线性空间，所以本章不建立新的微分理论，而是把已有的全微分、链式法则与Riesz表示用于矩阵，集中处理三件事：矩阵的形状，乘法不交换，以及对称、正定等结构约束。

第一节说明矩阵变量的微分与Frobenius梯度，并区分全部扰动与结构受限的扰动。后四节依次处理基本运算、求逆、行列式和对数行列式。每种运算都先确定定义域与输入输出，再由增量的一阶项读出微分并验证可微性；能写出精确增量的运算，就用同一个代数展开完成这两件事。

本章只建立微分公式与梯度。这些公式给出的一阶预测在有限扰动下有多准、误差被放大多少，留到第~\ref{chap:differential-perturbation-analysis}章统一估计。本章用到第~\ref{chap:linear-operations}章的迹、Frobenius内积、行列式与逆矩阵，第~\ref{chap:dual-spaces}章的Riesz表示，以及第~\ref{chap:solution-error-regularization}章逆矩阵的有限扰动界。

\section{矩阵微分与梯度}
\label{sec:multivariable-matrix-variable-gradient}

$m\times n$矩阵输入可以看成$mn$个实数坐标，$r\times s$矩阵输出可以看成$rs$个实数分量。按固定顺序（例如逐列）把矩阵元素排成长向量，矩阵空间就等同于欧氏坐标空间；第~\ref{sec:prelim-matrix-trace}节的Frobenius内积与范数$\langle\bm A,\bm B\rangle_F=\operatorname{tr}(\bm A^\top\bm B)=\sum_{i,j}A_{i,j}B_{i,j}$，正是拉直后向量的欧氏内积与欧氏范数。
本章的定义域$\mathcal U$都是全矩阵空间$\mathbb R^{m\times n}$中的开集，基点$\bm X\in\mathcal U$，因此全部充分小的同形扰动都可以使用。

设$\Phi:\mathcal U\to\mathbb R^{r\times s}$。若存在矩阵空间之间的线性映射$\mathrm d\Phi_{\bm X}:\mathbb R^{m\times n}\to\mathbb R^{r\times s}$，使
\[
 \Phi(\bm X+\bm H)=\Phi(\bm X)+\mathrm d\Phi_{\bm X}(\bm H)+o(\lVert\bm H\rVert_F),
\]
余项按输出的Frobenius范数度量，则称$\Phi$在$\bm X$处可微，$\mathrm d\Phi_{\bm X}$为其微分。这就是向量值可微性按矩阵形状重排后的写法，计算时常写$\mathrm d\Phi=\mathrm d\Phi_{\bm X}(\mathrm d\bm X)$。$\mathrm d\Phi_{\bm X}$是从输入扰动空间到输出空间的线性映射，拉直后它的矩阵是$rs\times mn$的Jacobian，一般不能写成一个与$\bm X$同形的矩阵。

还要分清两类“矩阵的函数”。$\bm X\mapsto\bm X^{-1}$、$\bm X\mapsto\bm X^2$这类映射按矩阵运算定义，输出的每个元素一般依赖输入的许多元素；逐元素作用（例如对每个元素分别取指数）则把各元素分别送入一元函数，微分只在对应位置上相乘。第~\ref{sec:ls-matrix-functions}节由标量函数沿谱构造的矩阵函数$f(\bm S)$属于前一类，但只是以矩阵为输入的映射中的一种；它与逐元素作用不能混为一谈，研究$\bm S\mapsto f(\bm S)$的微分时，扰动一般也不保持原来的特征方向。

\paragraph{输入变量决定扰动的形状}
同一个表达式随所选变量不同，微分的输入也不同。固定$\bm A$时，$\bm x\mapsto\bm A\bm x$是线性映射，微分为$\bm h\mapsto\bm A\bm h$，扰动是向量；固定$\bm x$时，$\bm A\mapsto\bm A\bm x$也是线性映射，微分为$\bm E\mapsto\bm E\bm x$，扰动是矩阵。两者同时变化时，把$(\bm A,\bm x)$看作乘积空间$\mathbb R^{m\times n}\times\mathbb R^n$中的一个输入，取乘积范数$\lVert(\bm E,\bm h)\rVert=(\lVert\bm E\rVert_F^2+\lVert\bm h\rVert_2^2)^{1/2}$。精确增量
\[
 (\bm A+\bm E)(\bm x+\bm h)-\bm A\bm x=\bm A\bm h+\bm E\bm x+\bm E\bm h
\]
中，$\lVert\bm E\bm h\rVert_2\leq\lVert\bm E\rVert_F\lVert\bm h\rVert_2\leq\tfrac12\lVert(\bm E,\bm h)\rVert^2$，所以联合微分为$(\bm E,\bm h)\mapsto\bm A\bm h+\bm E\bm x$，余项$\bm E\bm h$对联合扰动是二次量。

\paragraph{标量输出与矩阵梯度}
输出为标量（$r=s=1$）时，$\mathrm df_{\bm X}$是矩阵空间上的线性泛函。矩阵空间取Frobenius内积，由Riesz表示定理（定理~\ref{thm:ls-riesz-representation}），这个泛函可以唯一写成与一个同形矩阵的内积，这个矩阵就是梯度。这与向量情形$\mathrm df_{\bm a}(\bm h)=\nabla f(\bm a)^\top\bm h$完全对应。
\begin{definition}[矩阵标量函数的梯度]
\label{def:multivariable-matrix-scalar-gradient}
设$f:\mathcal U\to\mathbb R$在$\bm X$处可微。唯一满足
\begin{equation}
 \mathrm df_{\bm X}(\bm H)=\langle\nabla_{\bm X}f,\bm H\rangle_F
 =\operatorname{tr}\bigl((\nabla_{\bm X}f)^\top\bm H\bigr)
\label{eq:analysis-matrix-gradient-definition}
\end{equation}
的同形矩阵$\nabla_{\bm X}f\in\mathbb R^{m\times n}$称为矩阵梯度。
\end{definition}
取只有第$(i,j)$个元素为$1$的扰动$\bm H$，右侧内积只剩$\nabla_{\bm X}f$的第$(i,j)$个元素，左侧是对$X_{i,j}$的偏导数。所以矩阵梯度的各元素就是对相应矩阵元素的偏导数，唯一性也由此得到。与向量梯度一样，在Frobenius长度下，$\nabla_{\bm X}f$指出使$f$增加最快的扰动方向。

求梯度的一般做法是：先展开微分，再利用迹的循环移位把扰动移到末端，写成$\operatorname{tr}(\bm G^\top\mathrm d\bm X)$的形式，便读出$\nabla_{\bm X}f=\bm G$。乘积形状相容时，
\[
 \begin{gathered}
\operatorname{tr}(\bm A\bm B)=\operatorname{tr}(\bm B\bm A)\\
\operatorname{tr}(\bm A\bm B\bm C)=\operatorname{tr}(\bm B\bm C\bm A).
 \end{gathered}
\]
这里循环移动的是整个因子，不能任意交换相邻矩阵。标量等于自己的迹，所以任何标量表达式都可以放进迹中再移位。
\begin{example}[双线性函数的矩阵梯度]
\label{ex:multivariable-bilinear-matrix-gradient}
对固定向量$\bm a\in\mathbb R^m$、$\bm b\in\mathbb R^n$，令$f(\bm W)=\bm a^\top\bm W\bm b$。它对$\bm W$线性，微分为$\mathrm df=\bm a^\top(\mathrm d\bm W)\bm b$。把这个标量写成迹并循环移位，
\[
 \mathrm df=\operatorname{tr}\bigl(\bm a^\top(\mathrm d\bm W)\bm b\bigr)=\operatorname{tr}\bigl(\bm b\bm a^\top\mathrm d\bm W\bigr)
 =\operatorname{tr}\bigl((\bm a\bm b^\top)^\top\mathrm d\bm W\bigr),
\]
于是$\nabla_{\bm W}f=\bm a\bm b^\top$。逐元素核对：$f=\sum_{k,l}a_kW_{k,l}b_l$，对$W_{i,j}$的偏导数是$a_ib_j$，正是$\bm a\bm b^\top$的第$(i,j)$个元素。
\end{example}

\paragraph{结构化变量的梯度}
矩阵梯度的唯一性用到了全部矩阵方向。若变量限制在线性子空间$\mathcal S$中，例如对称矩阵空间，允许的扰动只能取$\mathcal S$中的方向。在$\mathcal S$上沿用Frobenius内积，Riesz表示定理在$\mathcal S$内给出唯一的表示矩阵，称为受限梯度。设$f$在全矩阵空间中有可微的延拓，环境梯度为$\bm G$，记$P_{\mathcal S}$为到$\mathcal S$的正交投影。$\bm G-P_{\mathcal S}\bm G$与$\mathcal S$正交（第~\ref{sec:ls-orthogonal-projection}节），所以对$\bm H\in\mathcal S$有$\langle\bm G,\bm H\rangle_F=\langle P_{\mathcal S}\bm G,\bm H\rangle_F$：受限梯度就是环境梯度的正交投影。

对称矩阵空间中，任一矩阵分解为对称部分与反对称部分，二者Frobenius正交，因为对称$\bm S$与反对称$\bm K$满足$\operatorname{tr}(\bm S\bm K)=\operatorname{tr}((\bm S\bm K)^\top)=-\operatorname{tr}(\bm K\bm S)=-\operatorname{tr}(\bm S\bm K)$；因此投影为$(\bm G+\bm G^\top)/2$。若在$\mathcal S$中仍使用环境梯度，满足$\mathrm df(\bm H)=\langle\cdot,\bm H\rangle_F$（$\bm H\in\mathcal S$）的矩阵就不唯一：$\bm G$加上任一与$\mathcal S$正交的矩阵都满足。

对称正定矩阵组成对称矩阵空间中的开集：由第~\ref{chap:condition-and-spectral-perturbation}章的Weyl界，对称扰动的谱范数小于最小特征值时正定性保持。所以讨论正定变量时，邻域和扰动方向都在对称子空间内理解。矩阵输出时没有这样的单个梯度矩阵，只能保留线性映射$\mathrm d\Phi_{\bm X}$本身。

\section{基本矩阵运算的微分}
\label{sec:multivariable-matrix-basic-operations}

本节从最简单的矩阵运算开始。线性运算的余项为零；乘法与平方的增量可以直接展开，微分之外剩下的部分就是可以精确写出的余项，由它验证可微性；迹、Frobenius范数平方与二次型则用来练习读出梯度。

\paragraph{线性运算}
加法、固定标量的数乘、转置与迹都是线性映射，精确增量就是它们作用于扰动的结果，余项为零，所以微分等于运算本身：
\[
 \begin{gathered}
 \mathrm d(\bm X+\bm Y)=\mathrm d\bm X+\mathrm d\bm Y,\\
 \mathrm d(c\bm X)=c\,\mathrm d\bm X,\\
 \mathrm d(\bm X^\top)=(\mathrm d\bm X)^\top,\\
 \mathrm d\operatorname{tr}\bm X=\operatorname{tr}(\mathrm d\bm X).
 \end{gathered}
\]
固定矩阵从两侧相乘也是线性的：$\mathrm d(\bm A\bm X\bm B)=\bm A(\mathrm d\bm X)\bm B$，扰动夹在中间的位置不能移动。复合运算沿用第~\ref{sec:analysis-vector-chain-matrix-calculus}节的链式法则：先计算内部扰动，再让外部微分作用于它。只要逐步核对每个线性映射的输入、输出形状，就不需要另背一套矩阵求导规则。

线性迹函数的微分已经是式~\eqref{eq:analysis-matrix-gradient-definition}的形式，可以直接读出梯度。对固定的$\bm{A}\in\mathbb{R}^{m\times n}$，令$f(\bm{X})=\operatorname{tr}(\bm{A}^{\top}\bm{X})=\sum_{i,j}A_{i,j}X_{i,j}$。由于$f$线性，$f(\bm X+\bm H)-f(\bm X)=\operatorname{tr}(\bm A^\top\bm H)$，没有余项，所以
\begin{equation}
  \nabla_{\bm{X}}\operatorname{tr}(\bm{A}^{\top}\bm{X})
  =\bm{A}.
\label{eq:analysis-trace-linear-gradient}
\end{equation}
这与向量情形$\nabla_{\bm x}(\bm a^\top\bm x)=\bm a$对应，常用于从迹表达式中读出梯度。

\paragraph{乘积的精确增量}
两个因子同时变化时，只用分配律就得到
\begin{equation}
 (\bm A+\bm E)(\bm B+\bm H)-\bm A\bm B=\bm A\bm H+\bm E\bm B+\bm E\bm H .
\label{eq:multivariable-product-increment}
\end{equation}
这一等式对任意大小的扰动都成立，本身不涉及微分。前两项对$(\bm E,\bm H)$线性，各被另一个因子放大：在满足乘积不等式的矩阵范数下，$\lVert\bm A\bm H+\bm E\bm B\rVert\leq\lVert\bm A\rVert\lVert\bm H\rVert+\lVert\bm E\rVert\lVert\bm B\rVert$。要把它们识别为微分，还须选定联合输入的大小，例如$\lVert(\bm E,\bm H)\rVert=(\lVert\bm E\rVert_F^2+\lVert\bm H\rVert_F^2)^{1/2}$。由Frobenius范数的次乘性，
\[
 \lVert\bm E\bm H\rVert_F\leq\lVert\bm E\rVert_F\lVert\bm H\rVert_F\leq\tfrac12\lVert(\bm E,\bm H)\rVert^2 ,
\]
交叉项除以联合扰动的大小后趋于零。所以
\[
 \mathrm D(\bm A\bm B)[\bm E,\bm H]=\bm A\bm H+\bm E\bm B,
\]
而$\bm E\bm H$是可以直接计算的精确余项。只固定一个因子时，乘积对另一个因子线性，交叉项不出现。

令两因子相同，$\Phi(\bm X)=\bm X^2$的精确增量为$(\bm X+\bm E)^2-\bm X^2=\bm X\bm E+\bm E\bm X+\bm E^2$。由$\lVert\bm E^2\rVert_F\leq\lVert\bm E\rVert_F^2$，得到$\mathrm d(\bm X^2)(\bm E)=\bm X\bm E+\bm E\bm X$，余项恰为$\bm E^2$。只有$\bm X$与扰动交换时，才能写成$2\bm X\bm E$。例如
\[
 \bm X=\begin{bmatrix}0&1\\0&0\end{bmatrix},\qquad
 \bm E=\begin{bmatrix}0&0\\1&0\end{bmatrix}
\]
时，$\bm X\bm E=\operatorname{diag}(1,0)$、$\bm E\bm X=\operatorname{diag}(0,1)$，正确的一阶项是$\bm I$，而照搬标量规则的$2\bm X\bm E=\operatorname{diag}(2,0)$是错的。直接核对：$(\bm X+t\bm E)^2=t\bm I$，对$t$的导数正是$\bm I$；本例中$\bm E^2=\bm 0$，余项为零。

\paragraph{Frobenius范数平方与二次型}
标量值运算的微分可以直接读出梯度。双线性形式$\bm a^\top\bm W\bm b$已在示例~\ref{ex:multivariable-bilinear-matrix-gradient}中处理。Frobenius范数平方的精确增量为
\[
 \lVert\bm X+\bm E\rVert_F^2-\lVert\bm X\rVert_F^2=2\langle\bm X,\bm E\rangle_F+\lVert\bm E\rVert_F^2 ,
\]
所以$\nabla_{\bm X}\lVert\bm X\rVert_F^2=2\bm X$，余项恰为$\lVert\bm E\rVert_F^2$。

同一个数组在固定线性作用和可变输入两种角色下，会产生不同的导数，二次型$\bm x^\top\bm A\bm x$是典型的例子。固定$\bm A$、以$\bm x$为变量时，由式~\eqref{eq:analysis-quadratic-form-gradient}，梯度为$(\bm A+\bm A^\top)\bm x$，是向量；固定$\bm x$、以$\bm A$为变量时，$\bm x^\top\bm E\bm x=\operatorname{tr}(\bm x\bm x^\top\bm E)=\langle\bm x\bm x^\top,\bm E\rangle_F$，梯度为$\bm x\bm x^\top$，是矩阵，且没有余项。两者同时变化时，展开得
\[
 \begin{aligned}
 &(\bm x+\bm h)^\top(\bm A+\bm E)(\bm x+\bm h)-\bm x^\top\bm A\bm x\\
 &\quad=\bm x^\top(\bm A+\bm A^\top)\bm h+\bm x^\top\bm E\bm x
 +\bigl(\bm h^\top\bm A\bm h+\bm x^\top\bm E\bm h+\bm h^\top\bm E\bm x+\bm h^\top\bm E\bm h\bigr).
 \end{aligned}
\]
括号中的每一项都至少含两个扰动因子，对联合扰动是高阶量，所以联合微分是两个单独微分之和：$(\bm E,\bm h)\mapsto\bm x^\top(\bm A+\bm A^\top)\bm h+\bm x^\top\bm E\bm x$。

\section{求逆的微分}
\label{sec:multivariable-inverse-differential}

求逆映射$\operatorname{inv}(\bm A)=\bm A^{-1}$定义在$d\times d$可逆矩阵组成的集合上，输入、输出都是方阵，属于矩阵输出的情形，没有单个梯度矩阵，要直接求微分。逆矩阵与行列式非零的关系已在第~\ref{sec:prelim-inverse-matrices}节及第~\ref{sec:prelim-determinants}节建立。本节的矩阵范数取由向量范数诱导的矩阵范数，它使$\lVert\bm I\rVert=1$并满足乘积不等式。

\paragraph{由逆差恒等式得到微分}
设$\bm A$可逆，$\bm B=\bm A+\bm E$，并设$q=\lVert\bm A^{-1}\rVert\lVert\bm E\rVert<1$。第~\ref{sec:ls-inverse-sensitivity}节已经证明，这时$\bm B$可逆，逆差满足精确恒等式~\eqref{eq:ch06-inverse-finite-perturbation}，即$\bm B^{-1}-\bm A^{-1}=-\bm B^{-1}\bm E\bm A^{-1}$，并且
\[
 \lVert\bm B^{-1}-\bm A^{-1}\rVert\leq\frac{\lVert\bm A^{-1}\rVert^2\lVert\bm E\rVert}{1-q}.
\]
特别地，可逆矩阵组成开集，而且$\bm E\to\bm0$时$\bm B^{-1}\to\bm A^{-1}$，求逆是连续的。

恒等式右端含未知的$\bm B^{-1}$，还不是对$\bm E$线性的式子。把其中的$\bm B^{-1}$换成$\bm A^{-1}$，就得到线性的候选项$-\bm A^{-1}\bm E\bm A^{-1}$。替换造成的差为
\[
 \bm B^{-1}-\bm A^{-1}+\bm A^{-1}\bm E\bm A^{-1}
 =-(\bm B^{-1}-\bm A^{-1})\bm E\bm A^{-1},
\]
其范数不超过$\lVert\bm B^{-1}-\bm A^{-1}\rVert\,\lVert\bm E\rVert\,\lVert\bm A^{-1}\rVert$。由上面的连续性，第一个因子随$\bm E\to\bm0$趋于零，所以这个差是$o(\lVert\bm E\rVert)$。有限维矩阵空间中各种范数等价，对Frobenius范数同样如此。因此求逆在每个可逆矩阵处可微，且
\begin{equation}
  \mathrm{d}(\bm{X}^{-1})
  =
  -\bm{X}^{-1}(\mathrm{d}\bm{X})\bm{X}^{-1}.
\label{eq:analysis-inverse-differential}
\end{equation}
这个导数以矩阵方向$\mathrm d\bm X$为输入，输出仍是同形的方阵；它是矩阵空间上的线性映射，不能写成单个梯度矩阵。上面的差还可以精确写出并给出二次界，这在第~\ref{sec:multivariable-inverse-perturbation}节用于估计一阶预测的误差。

这里先由精确恒等式确认了可微性。随后可以对恒等式$\bm{X}\bm{X}^{-1}=\bm{I}$两侧求微分作交叉核对：左侧用乘积法则，右侧是常矩阵，得到$(\mathrm{d}\bm{X})\bm{X}^{-1}+\bm{X}\,\mathrm{d}(\bm{X}^{-1})=\bm{0}$，左乘$\bm{X}^{-1}$便回到同一公式。这一形式推导以逆映射可微为前提，不能单独承担可微性的证明。

$1\times1$时这就是$(1/x)'=-1/x^2$。两个$\bm{X}^{-1}$位于扰动两侧，不能合并成$-\bm{X}^{-2}\mathrm{d}\bm{X}$，因为$\mathrm{d}\bm{X}$未必与$\bm{X}^{-1}$交换。
取一个对角矩阵和一个非对角扰动，就能看出两种写法的差别。
\begin{example}[逆矩阵微分中的乘法次序]
\label{ex:analysis-noncommuting-inverse-differential}
取$\bm X=\operatorname{diag}(2,1)$，扰动方向
\[
 \bm E=\begin{bmatrix}0&1\\1&0\end{bmatrix}.
\]
由式~\eqref{eq:analysis-inverse-differential}，沿$\bm X+t\bm E$求逆矩阵在$t=0$处的导数为
\[
 -\bm X^{-1}\bm E\bm X^{-1}
 =\begin{bmatrix}0&-1/2\\-1/2&0\end{bmatrix}.
\]
若照搬标量规则而写成$-\bm X^{-2}\bm E$，会得到
\[
 \begin{bmatrix}0&-1/4\\-1&0\end{bmatrix}.
\]
两种写法给出不同的矩阵。要判断哪一个正确，可以直接写出这条路径上的逆矩阵：
\[
 (\bm X+t\bm E)^{-1}=\frac1{2-t^2}\begin{bmatrix}1&-t\\-t&2\end{bmatrix}.
\]
在$t=0$处对$t$求导，两个非对角元的导数都是$-1/2$，对角元的导数为$0$，与第一种结果一致。
\end{example}

错误写法连对称性都丢失了。路径$\bm X+t\bm E$在$t=0$附近由对称可逆矩阵组成，它们的逆也对称，所以导数必须对称。这种结构检查能直接发现错误。

矩阵沿可微的参数曲线$\bm A(t)$变化，且每个$\bm A(t)$可逆时，链式法则与式~\eqref{eq:analysis-inverse-differential}给出
\[
 \frac{\mathrm d}{\mathrm dt}\bm A(t)^{-1}=-\bm A(t)^{-1}\bm A'(t)\bm A(t)^{-1}.
\]
第~\ref{sec:multivariable-logdet-differential}节和第~\ref{chap:differential-perturbation-analysis}章沿扰动线段求导时都会用到这个公式。

\section{行列式的微分}
\label{sec:multivariable-determinant-differential}

行列式把方阵映为一个实数，它的绝对值是线性映射对体积的缩放倍数（第~\ref{sec:prelim-determinants}节）。按排列展开，$\det\bm A$是各元素的多项式，所以它在全部方阵上连续可微，奇异矩阵处也不例外。本节由逐列展开算出微分，再在可逆时化为迹的形式。
\paragraph{代数余子式形式}
行列式对每一列分别线性（第~\ref{sec:prelim-determinants}节）。记$\bm A\in\mathbb R^{d\times d}$的列为$\bm a_1,\ldots,\bm a_d$，扰动$\bm E$的列为$\bm e_1,\ldots,\bm e_d$，逐列展开$\det(\bm a_1+\bm e_1,\ldots,\bm a_d+\bm e_d)$，得到$2^d$项，每一项在每一列取$\bm a_j$或$\bm e_j$。全取$\bm a_j$的一项是$\det\bm A$；恰有第$j$列取$\bm e_j$的一项，按第$j$列展开（第~\ref{sec:ls-det-cofactor-adjugate}节）等于$\sum_iE_{i,j}C_{i,j}$，其中$C_{i,j}$是$\bm A$的代数余子式，它不依赖$\bm A$的第$j$列。其余各项至少含两列扰动，是$\bm E$的元素的多项式，且每个单项至少二次；由$|E_{i,j}|\leq\lVert\bm E\rVert_F$，在$\lVert\bm E\rVert_F\leq1$时它们的和不超过某个常数乘$\lVert\bm E\rVert_F^2$。于是
\begin{equation}
 \mathrm d\det_{\bm A}(\bm E)=\sum_{i,j}C_{i,j}E_{i,j}=\langle\bm C,\bm E\rangle_F=\operatorname{tr}\bigl(\operatorname{adj}(\bm A)\bm E\bigr),
\label{eq:multivariable-determinant-cofactor-differential}
\end{equation}
其中$\bm C=[C_{i,j}]$是代数余子式矩阵，$\operatorname{adj}(\bm A)=\bm C^\top$是古典伴随。Frobenius梯度为$\nabla_{\bm A}\det=\bm C$。这个公式在奇异矩阵处同样成立。

$\bm A$可逆时，古典伴随恒等式（定理~\ref{thm:ls-adjugate-identity}）给出$\operatorname{adj}(\bm A)=\det(\bm A)\bm A^{-1}$，代入得到Jacobi公式的一阶形式
\[
 \mathrm{d}\det\bm{X}=\det(\bm{X})\operatorname{tr}(\bm{X}^{-1}\mathrm{d}\bm{X}),
\]
相应的梯度为$\det(\bm X)\bm X^{-\top}$。它说明有向体积缩放倍数的相对变化为
\[
 \frac{\mathrm d\det\bm X}{\det\bm X}=\operatorname{tr}(\bm X^{-1}\mathrm d\bm X).
\]
也可以用乘法公式核对：
\[
 \det(\bm X+\bm H)=\det(\bm X)\det(\bm I+\bm X^{-1}\bm H),
\]
问题化归为单位矩阵附近的一阶变化。奇异处没有$\bm X^{-1}$，只能使用代数余子式或古典伴随的形式。

在最简单的基点$\bm I$处，代数余子式矩阵就是$\bm I$，一阶变化是扰动的迹。
\begin{example}[单位矩阵附近行列式的一阶变化]
\label{ex:multivariable-determinant-near-identity}
取$\bm H=\operatorname{diag}(\varepsilon_1,\varepsilon_2)$，则
\[
 \det(\bm I+\bm H)=(1+\varepsilon_1)(1+\varepsilon_2)=1+\varepsilon_1+\varepsilon_2+\varepsilon_1\varepsilon_2,
\]
一阶变化是$\varepsilon_1+\varepsilon_2=\operatorname{tr}\bm H$。若扰动只在非对角位置，例如
\[
 \bm H=\begin{bmatrix}0&\varepsilon\\0&0\end{bmatrix},
\]
则$\det(\bm I+\bm H)=1$，面积不变，而$\operatorname{tr}\bm H=0$。
\end{example}

两个方向各拉长一点时，面积的相对增量一阶上是两个伸长率之和；剪切扰动不改变面积，迹也为零。两种情形中，一阶变化都等于扰动的迹。

\section{对数行列式的微分}
\label{sec:multivariable-logdet-differential}

实值的$\log\det\bm A$要求$\det\bm A>0$。行列式连续，满足这一条件的方阵组成开集，对数在其上与行列式复合。对行列式的微分应用标量对数的链式法则$\mathrm d\log u=\mathrm du/u$，得到
\begin{equation}
\begin{gathered}
\mathrm{d}\log\det\bm{X}
  =
  \operatorname{tr}(\bm{X}^{-1}\mathrm{d}\bm{X})\\
\nabla_{\bm{X}}\log\det\bm{X}
  =
  \bm{X}^{-\top}.
\end{gathered}
\label{eq:analysis-logdet-gradient}
\end{equation}
第二式由$\operatorname{tr}(\bm X^{-1}\mathrm d\bm X)=\operatorname{tr}\bigl((\bm X^{-\top})^\top\mathrm d\bm X\bigr)$对照式~\eqref{eq:analysis-matrix-gradient-definition}读出。
若$\bm X$还对称，则$\bm{X}^{-\top}=\bm{X}^{-1}$；一般情形应保留转置。$1\times1$时它回到$(\log x)'=1/x$。

\paragraph{对称正定变量}
统计与优化中常见的是对称正定矩阵上的$\log\det$。这时变量在对称矩阵子空间中，只允许对称扰动。由第~\ref{sec:multivariable-matrix-variable-gradient}节，受限梯度是环境梯度$\bm A^{-\top}=\bm A^{-1}$在对称子空间上的投影；它本身已对称，所以梯度就是$\bm A^{-1}$。

这里有两个不同的定义域。$\det>0$的方阵组成全矩阵空间中的开集，$\log\det$在其上已有定义，式~\eqref{eq:analysis-logdet-gradient}在这个开集上处处成立；对称正定矩阵组成对称子空间中的开集，是前者与对称子空间之交的一部分。所以正定性不是$\log\det$可微的条件，它只决定允许的扰动方向和梯度的表示。第~\ref{sec:multivariable-logdet-perturbation}节估计有限扰动下的余项时，正定性还提供整条扰动线段留在定义域内的明确保证。

\FloatBarrier
\section{本章小结}
\label{sec:matrix-differentiation-summary}

以矩阵为自变量时，微分仍是全部同形扰动上的线性映射，只是保留了矩阵形状。标量输出时，Riesz表示在Frobenius内积下给出同形的梯度矩阵，常用的读法是把微分写成$\operatorname{tr}(\bm G^\top\mathrm d\bm X)$；矩阵输出时没有单个梯度矩阵，只能保留微分本身，例如求逆的$\mathrm d\bm X\mapsto-\bm X^{-1}\mathrm d\bm X\bm X^{-1}$。变量限制在对称矩阵等子空间时，允许的方向随之受限，受限梯度是环境梯度的正交投影，对称情形为$(\bm G+\bm G^\top)/2$。

各运算的公式与条件可以这样归纳。线性运算的微分就是运算本身；乘积的微分是两个“只变一个因子”的项之和，因子顺序不能交换，平方的微分是$\bm X\mathrm d\bm X+\mathrm d\bm X\bm X$而不是$2\bm X\mathrm d\bm X$。求逆在可逆矩阵的开集上可微，扰动夹在两个逆矩阵之间。行列式是多项式，在包括奇异矩阵的全部方阵上可微，微分由代数余子式给出，可逆时化为$\det(\bm X)\operatorname{tr}(\bm X^{-1}\mathrm d\bm X)$。对数行列式在$\det>0$的开集上可微，梯度为$\bm X^{-\top}$，对称正定情形化为$\bm X^{-1}$。

所有公式都先由增量的一阶项识别、再由余项验证可微，对$\bm X\bm X^{-1}=\bm I$等恒等式形式求导只能作核对。精确增量中的余项本章只用来证明可微；这些公式作为一阶工具，在第~\ref{chap:differential-perturbation-analysis}章用来估计条件数和有限扰动下一阶预测的误差。
